\documentclass[UTF8,zihao=5,fontset=none]{ctexart}
\usepackage{clover}
\renewcommand{\doctitle}{项目说明}
% 单页：内容精简，行距与其他文档一致
\fancypagestyle{first}{\fancyhf{}}
\newgeometry{top=13mm,bottom=11mm,left=21mm,right=21mm}
\ctexset{section/beforeskip=7pt plus 1pt, section/afterskip=5pt}

\begin{document}

\cloverhead{Clover：投稿前的论文引用体检}{项目说明}{效率进化赛道\quad 阿里云特别赛题\quad Intel 本地 AI 应用挑战\quad 团队：我要发顶刊\quad 天津师范大学\quad 吕浩博}

\section{一句话介绍}

Clover 是给研究生和学者用的投稿前引用体检工具：拖入一篇 Word 或 LaTeX 稿件，几分钟内核对每一条参考文献是否真实存在、被引原文是否支持正文的论断、格式是否规范、引用是否符合目标期刊的常态，问题写成 Word 批注和修订交还作者逐条处理。

\section{用户需求}

投稿前逐条核对引用仍然靠人工，几十上百条参考文献很少有人真的逐条翻原文；大模型辅助写作又带来编造文献的新风险。错引在人工写作时代就普遍存在：Clover 在 10 篇 2016 至 2020 年发表的中英文论文里，报出并经人工核实的原文问题有 105 处。

\section{核心方案}

\begin{ctable}{@{}p{5.2em}L@{}}
\hd{检查} & \hd{做什么，结果怎么交给作者} \\ \midrule
文献真实性 & 到 Crossref、OpenAlex 等四个数据库逐条比对作者、年份、题名；查不到的中文文献标为无法核验 \\
论断支持度 & 在被引原文里检索相关片段，由 Qwen 判断是否支持，批注附逐字可查的证据句 \\
格式规范 & 错别字和参考文献重排写成 Word 修订，作者逐条接受或拒绝 \\
期刊对标 & 与目标期刊范文逐项比较章节结构、引用位置、参考文献构成，偏离常态处写批注 \\
原则 & 机械问题给修订，实质问题只给批注；不改写正文，不删除或替换参考文献 \\
\end{ctable}

\section{AI 能力}

\begin{itemize}
\item \textbf{云端：}阿里云百炼 Qwen3.8-flash 抽取论断、判断支持度、找错别字，Qwen3.8-max 复核。
\item \textbf{Intel 本机：}Qwen3-Embedding-0.6B 与 Qwen3-Reranker-0.6B 经 OpenVINO 跑在 Arc 核显或 NPU 上，挑出交给判断的原文段落；隐私模式下 Qwen3-8B 在核显上完成抽取和错别字检查。
\end{itemize}

\section{应用价值}

\begin{ctable}{@{}>{\raggedright\arraybackslash}p{12em}L@{}}
\hd{指标} & \hd{结果} \\ \midrule
人为注入的引用错误 & 10 篇中英文稿件注入 86 处，检出 84 处 \\
报出的问题里属实的比例 & 公开数据集 Citation-Integrity 上为 73\%，开源模型 SemanticCite 为 50\% \\
单篇云端花费与耗时 & 全云端 2.43 元；向量与重排放到 Intel 本机后 0.40 元，约 9 分钟 \\
隐私模式发往云端的文字 & 11,700 字中 5,910 字，全部是带引用的论断句 \\
\end{ctable}


\end{document}
